\documentclass[10pt,twocolumn,letterpaper]{article}

\usepackage{booktabs}
\usepackage{graphicx}
\usepackage{amsmath,amssymb}
\usepackage{microtype}
\usepackage{xcolor}
\usepackage[hidelinks]{hyperref}

\newcommand{\method}{\textsc{LaFGS}}
\newcommand{\R}{\mathbb{R}}
\newcommand{\Gprior}{\mathcal{G}_{\mathrm{rgb}}}
\newcommand{\Tracks}{\mathcal{T}}
\newcommand{\Anchors}{\mathcal{A}}
\newcommand{\Queries}{\mathcal{Q}}
\newcommand{\TODO}[1]{\textcolor{red}{[TODO: #1]}}


\title{LaFGS: Rendering-to-Localization Topology Distillation\\
for Compact Gaussian-Supported Visual Maps}
\author{Anonymous Authors}

\begin{document}
\maketitle

\begin{abstract}
Photometric Gaussian primitives are designed to explain image formation, not
to serve as sparse localization landmarks. We present \method, an offline map
reconstruction method that distills an arbitrary frozen RGB Gaussian prior
into a compact track-centric localization layer. Real mapping images define
cross-view feature identities and independently triangulated geometry, while
Gaussian raster provenance supplies surface support without imposing a
one-to-one primitive-to-landmark correspondence. The map then repeatedly
localizes its mapping views and uses the resulting retrieval outcomes to
reconstruct descriptors and select a pose-sufficient topology. At deployment,
the learned map requires only native SuperPoint extraction, global cosine
top-1 matching, and one standard PnP/RANSAC solve. Experiments are designed to
test localization accuracy, map compactness, prior flexibility, and transfer
across Cambridge Landmarks, 7Scenes, and 12Scenes.
\end{abstract}

\section{Introduction}

Gaussian Splatting reconstructs an explicit scene representation that is
effective for rendering. Its primitives, however, inherit their number,
support, and geometry from a photometric objective. A single surface region can
be represented by many splats; one splat can support several repeatable image
features; and floaters that are tolerable for rendering can become harmful 3D
correspondences. Consequently, a rendering primitive is not necessarily a
localization landmark.

We study a different objective: reconstruct a sparse localization layer from
a frozen, off-the-shelf RGB Gaussian scene prior. The desired map should retain
landmarks that are repeatedly detected, descriptor-distinctive, geometrically
reliable, useful to PnP, and compact. Importantly, the expensive reasoning is
restricted to offline map construction. Online localization remains the
standard sparse pipeline of descriptor extraction, one global top-1 match per
keypoint, and one PnP/RANSAC solve.

\method addresses this problem through \emph{rendering-to-localization topology
distillation}. Real mapping images first define track identities independently
of Gaussian IDs. Robust triangulation supplies PnP geometry, and raster
provenance connects each track-centric anchor to the supporting Gaussian
surface. A self-localization loop then applies the current map to the mapping
views and learns from kept correspondences, swaps, missed positives, and false
attractors. Finally, a pose-coverage-aware distillation stage produces a compact
single-descriptor map.

Our contributions are:
\begin{itemize}
  \item Gaussian-supported track-centric anchors that separate rendering
        primitive identity, feature-track identity, and localization identity.
  \item Self-localization-guided descriptor reconstruction that learns from the
        current map's sparse retrieval outcomes rather than static feature
        averaging alone.
  \item Task-specific topology distillation that converts a large candidate
        universe into a compact, pose-sufficient localization map.
  \item A prior-flexible construction protocol that keeps the RGB Gaussian
        frozen and reconstructs localization geometry from real image tracks,
        limiting primitive-level error propagation.
\end{itemize}

\section{Related Work}

\paragraph{Sparse visual localization.}
Classical structure-based localization matches local image descriptors to a 3D
landmark map and estimates pose with a robust geometric solver. Learned local
features improve repeatability and matching, but the landmark topology is
usually inherited from SfM or selected after reconstruction. \method retains
the sparse deployment interface while moving task-specific learning into map
construction.

\paragraph{Gaussian scene representations.}
3D and 2D Gaussian Splatting optimize explicit primitives for photometric
rendering. Several localization systems render appearance or attach features
to Gaussian primitives. Our formulation instead treats the Gaussian model as a
frozen geometric and visibility scaffold: localization identities and PnP
points are reconstructed from real feature tracks and need not correspond
one-to-one with primitives.

\paragraph{Compact localization maps.}
Map compression commonly ranks or prunes existing 3D points according to
visibility, descriptor quality, or pose coverage. Our topology is reconstructed
jointly with descriptor evidence: track-derived anchors form the core, Gaussian
lineage supplies coverage candidates, and mapping-view self-localization
provides task-specific functional evidence.

\TODO{Replace this concise positioning with citation-grounded discussion before
submission.}

\section{Method}

\subsection{Problem Formulation}

Let a frozen RGB Gaussian prior be
$\Gprior=\{(\boldsymbol\mu_j,\boldsymbol\Sigma_j,
\alpha_j,\mathbf{c}_j)\}_{j=1}^{N_G}$ and let the mapping set
$\Queries=\{(I_q,K_q,T_q)\}$ contain RGB images with calibrated poses. We
reconstruct a localization map
\begin{equation}
  \Anchors=\{(\mathbf{x}_i,\mathbf{d}_i,\ell_i)\}_{i=1}^{N_A},
\end{equation}
where $\mathbf{x}_i\in\R^3$ is PnP geometry, $\mathbf{d}_i$ is one deployment
descriptor, and $\ell_i$ stores Gaussian support lineage. In general,
\begin{equation}
  \text{Gaussian ID}\neq\text{track ID}\neq\text{anchor ID}.
\end{equation}
The prior remains frozen throughout localization-map reconstruction.

\subsection{KCS/GWFF Bootstrap}

Native SuperPoint observations are extracted from every mapping image. We
select a broad Gaussian-supported scaffold with keypoint-consensus sampling
(KCS): a primitive must receive support from multiple views, viewpoint bins,
and trajectory bins. Geometry-weighted feature fusion (GWFF) initializes each
selected descriptor from visible observations, weighting incidence and view
geometry and trimming low-consensus descriptor observations. This bootstrap is
an initialization, not the final landmark identity.

\subsection{Track-First Evidence and Geometry}

We build reciprocal, cycle-consistent cross-view associations from the native
observations. A track $t\in\Tracks$ is accepted only when its observations span
sufficient views and viewpoint bins. Its 3D position is estimated independently
of Gaussian centers by robust multi-view triangulation:
\begin{equation}
  \mathbf{x}_t^*=
  \arg\min_{\mathbf{x}}\sum_{(q,\mathbf{u})\in t}
  \rho\!\left(\left\|\pi(K_q,T_q,\mathbf{x})-\mathbf{u}\right\|_2\right),
\end{equation}
subject to parallax, reprojection, and covariance gates. Gaussian rasterization
then returns visible primitive IDs and composition mass at each observation.
Aggregating this provenance gives surface support $\ell_t$ without replacing
$\mathbf{x}_t^*$ by a primitive center.

\subsection{Self-Localization-Guided Reconstruction}

At iteration $s$, the current map retrieves a top-$K$ candidate set for each
native mapping keypoint. Known mapping poses label outcomes relative to the
current map. The reconstruction objective separates:
\begin{equation}
  \mathcal{L}_{\mathrm{SL}} =
  \lambda_k\mathcal{L}_{\mathrm{keep}}+
  \lambda_s\mathcal{L}_{\mathrm{swap}}+
  \lambda_m\mathcal{L}_{\mathrm{miss}}+
  \lambda_a\mathcal{L}_{\mathrm{attractor}}+
  \lambda_l\mathcal{L}_{\mathrm{local}}+
  \lambda_t\mathcal{L}_{\mathrm{trust}}.
\end{equation}
The keep term protects already correct high-margin assignments. Swap and miss
terms rank a geometrically valid positive above the selected wrong candidate or
retrieve a positive absent from the current top ranks. The global-attractor
term suppresses landmarks that collect false matches across queries. A weak
local peak objective sharpens spatially nearby alternatives, while a protected
set prevents such updates from damaging pose-critical correspondences. Because
candidates are regenerated from the current descriptors, this is an online
matching loop during offline reconstruction rather than static post-processing.

\subsection{Localization Topology Distillation}

The candidate universe contains track-centric anchors and Gaussian-supported
coverage candidates. We first select a reliable track core using cross-view and
geometric evidence. A coverage reserve then adds candidates only where mapping
queries lack enough usable rows, with pose-information scoring preventing a
purely appearance-driven map. The resulting topology has a scene-dependent
size determined by the frozen selection rule rather than a manually chosen
per-scene target.

On this fixed compact topology, a short low-rank metric update uses the complete
positive teacher and group-robust weighting. The deployed descriptor remains a
single vector per anchor; no family prototypes, learned sampler, dense stage,
or test-time rendering is required.

\subsection{Sparse Deployment}

For a query image, native SuperPoint produces keypoints and descriptors. A
shared metric transforms query and map descriptors, cosine top-1 matching
produces one correspondence per keypoint, and one standard PoseLib PnP/RANSAC
call estimates the pose. All correspondences are passed to the solver without a
landmark-side match cap or a learned test-time selector.

\section{Experiments}

\subsection{Protocol}

We freeze one method configuration across scenes. Mapping images are used for
Gaussian-prior construction and all offline localization-map learning; official
test images are used only for evaluation. The deployment frontend uses native
SuperPoint, full-resolution coordinates with the explicit pixel-center
convention, global cosine top-1 matching, and one PoseLib PnP/RANSAC solve.

We report median, mean, and 90th-percentile translation and rotation error;
recall at $5\,$cm/$5^\circ$; raw and RANSAC-inlier correspondence precision at
2 and 4 pixels; catastrophic failures; primitive, track, and final anchor
counts; map memory; build time; and localization P50/P90 latency.

\subsection{Datasets and Baselines}

The evaluation covers five Cambridge Landmarks scenes, all seven 7Scenes
scenes, and all twelve 12Scenes scenes. The indoor main tables use the same
off-the-shelf Gaussian-prior method without semantic masks. A representative
prior-flexibility subset compares vanilla 3DGS, vanilla 2DGS, an enhanced 2DGS
prior, and a feed-forward AnySplat prior while keeping every localization-map
stage fixed.

Our principal comparison is between the Gaussian-supported KCS/GWFF bootstrap
(A0) and the complete reconstructed map (A1). Track-only and base-only maps are
reported on representative scenes to isolate track identity and Gaussian
coverage.

\begin{table*}[t]
  \centering
  \caption{Frozen-protocol localization results. Values are populated only by
  the registered experiment summaries.}
  \label{tab:main-results}
  \resizebox{\textwidth}{!}{%
  \begin{tabular}{llrrrrrrrr}
    \toprule
    Dataset & Method & Med. TE & Mean TE & P90 TE & Med. AE & 5/5 & Raw P@2 & Anchors & P90 time \\
    \midrule
    Cambridge & A0 & \multicolumn{8}{c}{\TODO{aggregate five-scene results}} \\
    Cambridge & A1 & \multicolumn{8}{c}{\TODO{aggregate five-scene results}} \\
    7Scenes & A0 & \multicolumn{8}{c}{\TODO{run seven official scenes}} \\
    7Scenes & A1 & \multicolumn{8}{c}{\TODO{run seven official scenes}} \\
    12Scenes & A0 & \multicolumn{8}{c}{\TODO{run twelve official scenes}} \\
    12Scenes & A1 & \multicolumn{8}{c}{\TODO{run twelve official scenes}} \\
    \bottomrule
  \end{tabular}}
\end{table*}


\subsection{Analysis Plan}

The evaluation tests four questions: whether A1 consistently improves A0;
whether compact topology preserves both central and tail pose accuracy; whether
the reconstructed map is robust to off-the-shelf prior quality; and whether the
offline cost yields a simpler online pipeline. A scene failure is attributed
only after auditing dataset convention, prior reconstruction, track coverage,
and candidate precision under the frozen protocol.

\section{Conclusion}

We formulate localization-map construction from Gaussian scenes as
rendering-to-localization topology distillation. The key separation between
Gaussian primitives, image tracks, and localization anchors allows real
cross-view evidence to define identity and PnP geometry while retaining the
Gaussian prior as a useful visibility and surface scaffold. Self-localization
reconstructs descriptors offline, and topology distillation produces a compact
map that deploys with a conventional sparse solver. The completed experiments
will establish the method's cross-domain accuracy, compactness, prior
flexibility, and computational trade-offs.


\bibliographystyle{plain}
\bibliography{references}
\end{document}
